\section{0x12\ 基础\ 算法}
\label{sec:0x12}

\subsection{主定理}

\Highlight{主定理}：用于求解分治递归类型的渐进时间复杂度问题。

具体地，对于一个规模为 $n$ 的问题，将其划分成 $a$ 个规模为 $\frac{n}{b}$ 的子问题，附带 “合并/分割” 的开销 $f(n)$。
\[
T(n) = aT\left( \frac{n}{b} \right) + f(n)
\]

\textbf{以 $n^{\log_b a}$ 为处理问题用时的基准线。}

\begin{itemize}
\item 若 $f(n) = O(n^{\log_b a - \epsilon})$\newline{}
\textbf{直觉：总工作量由 “叶子层” 主导。}
\[
T(n) = \Theta(n^{\log_b a})
\]

\item 若 $f(n) = \Theta(n^{\log_b a}\log^k n)$ 且满足 $k\geq 0$\newline{}
\textbf{直觉：每一层工作量相近，共 $\log_b n$ 层。}
\[
T(n) = \Theta(n^{\log_b a}\log^{k + 1} n)
\]

\item 若 $f(n) = \Omega(n^{\log_b a + \epsilon})$ 且满足正则性条件 $af(\frac{n}{b}) \leq cf(n)$（其中 $c < 1$）\newline{}
\textbf{直觉：总工作量由顶层的 合并/分割 主导。}
\[
T(n) = \Theta(f(n))
\]
\end{itemize}

\begin{quoteblock}
一些例子：

\begin{itemize}
\item 情形 1：遍历满二叉树 $T(n) = 2T(\frac{n}{2}) + O(1) = O(n)$。
\item 情形 2：归并排序 $T(n) = 2T(\frac{n}{2}) + O(n) = \mathcal{O}(n \log n)$。
\item 情形 3：多项式牛顿迭代 $T(n) = T(\frac{n}{2}) + \mathcal{O}(n \log n) = \mathcal{O}(n \log n)$。
\end{itemize}

一个主定理处理不了的例子：$T(n) = 2T\left(\frac{n}{2}\right) + \mathcal{O}\left(\frac{n}{\log n}\right) = \mathcal{O}(n \log \log n)$。
\end{quoteblock}

\subsection{快速乘}

求解 $a \times b \bmod p$，其中 $a, b, p$ 均为 \texttt{long long} 范围。

正道正确率更高，魔道效率更高。

\subsubsection{正道：i128}

\InputCode{cpp}{contents/0x10 基础/code/0x12 快速乘（i128）.cpp}

\subsubsection{魔道：long double}

注意到 $a\times b \bmod p = a\times b - \lfloor \frac{a \times b}{p} \rfloor\times p$，利用 \texttt{long double} 来处理 $\lfloor \frac{a \times b}{p} \rfloor$。

\textbf{模数较大时可能会出锅。}

\InputCode{cpp}{contents/0x10 基础/code/0x12 快速乘（long double）.cpp}

\subsection{光速幂}

对于一个\textbf{固定底数} $a$，多次询问 $a^i \bmod p$ 的值（其中 $0 \leq i \leq n$）。

取阈值 $B = \left\lceil \sqrt{n} \right\rceil$，先预处理出 $a^1, a^2, \dots, a^i, \dots, a^B$，再预处理出 $a^B, a^{2B}, \dots, a^{iB}, \dots, a^{\lfloor n / B \rfloor B}$。

每次只需利用 $a^i = a^{\lfloor i / B \rfloor B}\times a^{i \bmod B}$ 回答询问即可。

时间复杂度：预处理 $\mathcal{O}(\sqrt{n})$，查询 $\mathcal{O}(1)$。

\begin{quoteblock}
$k$ 进制快速幂：利用 $a^i = \left( a^k \right)^{\lfloor i / k \rfloor} \times a^{i \bmod k}$ 递归处理询问（右侧预处理，左侧递归）。

时间复杂度：预处理 $\mathcal{O}(k \log_k n)$，查询 $\mathcal{O}(\log_k n)$。
\end{quoteblock}

\InputCode{cpp}{contents/0x10 基础/code/0x12 光速幂.cpp}

